\begin{lemma}
\label{lemma-characterize-formally-unramified}
Let $R \to S$ be a ring map.
The following are equivalent:
\begin{enumerate}
\item $R \to S$ is formally unramified,
\item the module of differentials $\Omega_{S/R}$ is zero.
\end{enumerate}
\end{lemma}

\begin{proof}
Let $J = \Ker(S \otimes_R S \to S)$ be the kernel of
the multiplication map. Let $A_{univ} = S \otimes_R S/J^2$. Recall
that $I_{univ} = J/J^2$ is isomorphic to $\Omega_{S/R}$, see
Lemma \ref{lemma-differentials-diagonal}. Moreover, the two $R$-algebra maps
$\sigma_1, \sigma_2 : S \to A_{univ}$, $\sigma_1(s) = s \otimes 1 \bmod J^2$,
and $\sigma_2(s) = 1 \otimes s \bmod J^2$ differ by the
universal derivation $\text{d} : S \to \Omega_{S/R} = I_{univ}$.

\medskip\noindent
Assume $R \to S$ formally unramified.
Then we see that $\sigma_1 = \sigma_2$.
Hence $\text{d}(s) = 0$ for all $s \in S$.
Hence $\Omega_{S/R} = 0$.

\medskip\noindent
Assume that $\Omega_{S/R} = 0$. Let $A, I, R \to A, S \to A/I$
be a solid diagram as in Definition \ref{definition-formally-unramified}.
Let $\tau_1, \tau_2 : S \to A$ be two dotted arrows making the
diagram commute. Consider the $R$-algebra map $A_{univ} \to A$
defined by the rule $s_1 \otimes s_2 \mapsto \tau_1(s_1)\tau_2(s_2)$.
We omit the verification that this is well defined. Since $A_{univ} \cong S$
as $I_{univ} = \Omega_{S/R} = 0$ we conclude that $\tau_1 = \tau_2$.
\end{proof}
